% Nouveau EGA V (cote 25): archive cote 25, site page 52 (PDF page 53).
% typescript. Draft compared: none.
% AI-transcribed by Claude Opus 5.5 (claude-opus-5-5, adaptive thinking, effort high), Message Batches API, batch ega25b (msgbatch_01A7n5YuXubUPNK1zsepi8RG), 2026-11-13; not yet checked by eye.
\scanpage{25/52}
\mspage{\hand{47}}

\var{\xout{En vertu de 5.6, l'ensemble U des P tels que Y soit géom}}{}{typewriter strike-out}{ega25_25_052-v0001}

Supposons d'abord X géom normale. On sait déjà par 8.14. a)  que
l'en \var{\xout{ex}}{}{typewriter strike-out}{ega25_25_052-v0002} peut trouver un fermé Z' de P de codi\unsure{mension} $\underset{\hand{\rule{0.7em}{0.35pt}}}{>}$2, tel que
\hand{$\xi\in$} P - Z'  implique que  $Y_{\hand{\xi}}$  est  \uline{séparable} sur k(\hand{$\xi$}) de  dim \var{\inserted{d-}}{d-}{insertion}{ega25_25_052-v0003}1; pour
un tel \hand{$\xi$} , il revient donc au même que $Y_{\hand{\xi}}$ soit géom irréd, ou géom
intègre, et  les deux problèmes  (respé et non respé) examinés dans
8.18. sont donc équivalents\unsure{.|!} D'autre part, en vertu de 5.6. l'ensemble
\var{\inserted[typed]{U}}{U}{insertion}{ega25_25_052-v0004} des \hand{$\xi\in$} $\overset{\hand{\vee}}{P}$  tels que  $Y_{\hand{\xi}}$.  soit géom intègre de dim \var{\over{1}{d}}{d}{writing over}{ega25_25_052-v0005}-1 (l'hypothèse
dimensionnelle exprimant que $\hand{\phi}_{\hand{\xi}}$ est $X_{k(\hand{\xi})}$-régulier ) est \uline{ouvert}.
Nous allons exhiber un ensemble ouvert \var{\inserted[typed]{non vide,}}{non vide,}{insertion}{ega25_25_052-v0006} évident  P - Z  cont\unsure{è|e}nu dans U,
en\var{\xout{p}}{}{typewriter strike-out}{ega25_25_052-v0007} prenant pour  Z la réunion de  $\overset{\hand{\rule{3.3em}{0.35pt}}}{g(\underline{Y}'^{\,\mathrm{sing}})}$  et des hyperplans $\overset{\hand{\vee}}{H}_x$
de $\overset{\hand{\vee}}{P}$  définis par les  f(x), \hand{$x\in$} T. \var{\xout{Pour}}{}{typewriter strike-out}{ega25_25_052-v0008} Pour \hand{$\xi\in$} $\overset{\hand{\vee}}{P}$ - Z,  $Y_{\hand{\xi}}$  est \uline{lisse}
de dimension \var{\inserted{d-}}{d-}{insertion}{ega25_25_052-v0009}1, et comme il est géom connexe par 6.1. , il est géom
intègre.  Il faut donc  exprimer qu\unsure{e|è}\var{\xout{chaque}}{}{typewriter strike-out}{ega25_25_052-v0010} \var{\inserted[typed]{toute}}{toute}{insertion}{ega25_25_052-v0011} composante irréductible
de \var{\inserted[typed]{co}}{co}{insertion}{ega25_25_052-v0012}dimension  1 de Z  \var{\xout{ne ren}}{}{typewriter strike-out}{ega25_25_052-v0013}\var{\inserted{ren}}{ren}{insertion}{ega25_25_052-v0014}contre l'ouvert U.  Or ces composantes
irréductibles sont les  $\overset{\hand{\vee}}{H}_x$  (répetés éventuellement, peu importe), et
d'autre part $\overset{\hand{\rule{3.3em}{0.35pt}}}{g(\underline{Y}'^{\,\mathrm{sing}})}$  lorsque  ce dernier  est bien de codim 1, i.e.
X' ``non réglée'' pour f, (NB on utilise l'irréductibilité de  $\underline{Y}'^{\,\mathrm{sing}}$ ).
D'ailleurs, pour que ce dernier ensemble rencontre  l'ouvert U,
il f et s  que  g($\underline{Y}'^{\,\mathrm{sing}}$)  (qui en contient un ouvert dense) rencontre
U, \var{\xout{i.e. qu'il existe x' X'(k) et une sect}}{}{typewriter strike-out}{ega25_25_052-v0015} Cela prouve 8.18. dans ce
cas.  Lorsqu'on ne suppose pas X normal, on applique le résultat précé-%
dent au normalisé de X; le raisonnement est immédiat et je ne le détaille
pas ici.  \var{\margin{\uline{NB} dans le cas \var{\inserted{\unsure{bis|bas}}}{\unsure{bis|bas}}{insertion}{ega25_25_052-v0017}respé, 8.18. est contenu dans 8.16. plus précis sauf pour d=2. C'est pour le cas non respé qu'il peut être utile de ne pas exiger d=2, mais seulement $d\geq 2$ ....}}{\placedmargin{\uline{NB} dans le cas \var{\inserted{\unsure{bis|bas}}}{\unsure{bis|bas}}{insertion}{ega25_25_052-v0017}respé, 8.18. est contenu dans 8.16. plus précis sauf pour d=2. C'est pour le cas non respé qu'il peut être utile de ne pas exiger d=2, mais seulement $d\geq 2$ ....}}{marginal note placed}{ega25_25_052-v0016}

Il reste à expliciter  les conditions a) et b) de 8.18. Cela con-%
duit à examiner de façon générale la situation suivante:  On suppose
X \var{\deleted{intègre}}{}{deletion}{ega25_25_052-v0018}\var{\inserted{géom irréd,}}{géom irréd,}{insertion}{ega25_25_052-v0019} sur k \var{\deleted{alg clos}}{}{deletion}{ega25_25_052-v0020}, on se donne une sous-variété linéaire  L \hand{de}

\var{\margin{\hand{52}}}{\placedmargin{\hand{52}}}{marginal note placed}{ega25_25_052-v0021}
